\documentclass{article}

% Use numeric, compressed citations.
\PassOptionsToPackage{numbers, compress}{natbib}

% Show author names + a "work in progress" footer (this is a progress report).
% For the anonymized submission build, switch to: \usepackage{neurips_2023}
\usepackage[preprint]{neurips_2023}

\usepackage[utf8]{inputenc} % allow utf-8 input
\usepackage[T1]{fontenc}    % use 8-bit T1 fonts
% For Hangul author names, compile with XeLaTeX/LuaLaTeX and uncomment:
% \usepackage{kotex}
\usepackage{hyperref}       % hyperlinks
\usepackage{url}            % simple URL typesetting
\usepackage{booktabs}       % professional-quality tables
\usepackage{amsfonts}       % blackboard math symbols
\usepackage{amsmath}        % align, etc.
\usepackage{nicefrac}       % compact symbols for 1/2, etc.
\usepackage{microtype}      % microtypography
\usepackage{xcolor}         % colors


\title{Multi-Agent Reinforcement Learning for Mario Kart Wii:\\
Rank-Conditioned Item Strategy and Team Cooperation under Rubber-Banding}


\author{%
  Kangmin Kim (20250064) \quad Dongha Kwak (20250041) \quad Seojin Yoon (20250482) \\
  Team 3, CS377 \\
  KAIST \\
  \texttt{\{kangmin.kim, dongha.kwak, seojin.yoon\}@kaist.ac.kr} \\
}


\begin{document}


\maketitle


\begin{abstract}
Mario Kart Wii (MKW) couples a \emph{rubber-banding} item mechanic---trailing
players receive stronger items than leaders---with team racing, producing
asymmetric, state-dependent dynamics under a shared, rank-based reward. We study
a 2v2 setting (four karts, two teams) controlled by a single shared policy, and
ask two questions: (i) does the shared policy learn \emph{rank-differentiated}
behavior (defensive item-saving at rank~1 vs.\ aggressive immediate use at
rank~4), or collapse to a single mode; and (ii) how does intra-team cooperation
emerge. Building on the Dolphin emulator and memory-reading infrastructure of
Beyond the Rainbow (BTR), we depart from it in three ways: multi-agent self-play
instead of single-agent play, structured memory-state observations instead of
images, and MAPPO with a centralized critic instead of Rainbow DQN. This report
presents the full method and a verified end-to-end implementation: an
information-asymmetric observation design (full self/teammate, coarse
opponents), a potential-based rank reward balancing individual rank with
team success, a centralized-critic MAPPO learner, and a frozen-snapshot
self-play loop. We report environment-verification findings (notably a corrected
item-button mapping without which items can never be used), unit and integration
test results, two methodological refinements developed during training---an
identical-kart control that isolates rank from vehicle-stat confounds, and a
race-first curriculum that learns driving before item dynamics---and the planned
experimental protocol for the two research questions.
\end{abstract}


\section{Introduction}

Mario Kart Wii's item system is explicitly \emph{rank-conditioned}: a kart in
last place is far more likely to draw powerful catch-up items (e.g.\ the
Bullet Bill, lightning) than the leader, who tends to draw defensive or weak
items. This rubber-banding makes the value of an item, and therefore the optimal
\emph{timing} of its use, depend on the agent's current rank. Intuitively, a
leading agent should hoard a defensive item (e.g.\ trail a banana to block an
incoming shell), whereas a trailing agent should use offensive items
immediately to climb. Whether a single, rank-agnostic \emph{shared} policy can
recover this dichotomy purely from a rank-based reward---rather than collapsing
to one mode---is our first research question (RQ1). In the 2v2 team setting, a
second question (RQ2) concerns the emergence of cooperation: protecting a
teammate and coordinating attacks against opponents.

We adopt the environment and memory-reading toolchain of Beyond the Rainbow
(BTR)~\citep{btr}, an image-based Rainbow~DQN agent that reaches high placements
against fixed CPU opponents in single-agent MKW. We depart from BTR in three
ways: (1) we move from single-agent play to multi-agent self-play; (2) we replace
image inputs with structured memory-state observations (position, velocity,
rank, lap, item, boost state); and (3) we replace Rainbow DQN with MAPPO, a
centralized-critic actor--critic method for cooperative multi-agent
games~\citep{mappo}. Training runs in VS Race mode on a fixed track with items
enabled, under centralized training with decentralized execution (CTDE).

This document is a progress report. The complete system has been implemented and
verified end-to-end; the empirical study of RQ1 and RQ2 is the remaining work
and its protocol is specified in Section~\ref{sec:plan}.


\section{Background and Related Work}

\paragraph{MARL and CTDE.} We follow the centralized-training,
decentralized-execution paradigm~\citep{maddpg}: actors condition only on local,
restricted observations available at execution, while a centralized critic may
access global state during training. MAPPO~\citep{mappo} instantiates CTDE with
PPO~\citep{ppo} actors and a centralized value function, and is a strong,
simple baseline for cooperative control; we use a single \emph{shared} actor
across all four karts, with a team indicator encoded in the observation.

\paragraph{Potential-based reward shaping.} To turn the sparse rank signal into a
dense, stable learning target without distorting the optimal policy, our rank
rewards are \emph{potential-based}~\citep{shaping}: a reward of the form
$\gamma\,\Phi(s')-\Phi(s)$ leaves the set of optimal policies unchanged while
densely rewarding rank gains.

\paragraph{Self-play.} To maintain opponent diversity and reduce strategy
cycling, we train the learning team against frozen past snapshots of the shared
policy, a standard competitive self-play technique~\citep{selfplay}.


\section{Method}

\subsection{Environment and observations}

The environment wraps the Dolphin emulator as a four-player local race and
exposes per-kart state read directly from game memory each step
(\texttt{race\_position}, \texttt{current/max\_race\_completion},
\texttt{current\_lap}, \texttt{speed}, position/velocity/rotation,
\texttt{drift\_state}, \texttt{item}, mini-turbo charge and boost timers).
Teams are fixed as agents $\{0,1\}$ vs.\ $\{2,3\}$.

Following the proposal, each actor receives a \emph{restricted} view modeling
realistic perception with perfect intra-team communication: full state of itself
and its teammate, but only \emph{coarse} information about opponents (relative
position and rank---what a human reads from the screen and minimap). The
centralized critic instead receives the full global state of all four karts.
Table~\ref{tab:obs} summarizes the decomposition; concretely the actor
observation is $59$-dimensional and the critic state is $102$-dimensional.

\begin{table}[t]
  \caption{Information-asymmetric observation design (CTDE). The actor never
  observes opponent item/velocity/charge; the centralized critic observes
  everything.}
  \label{tab:obs}
  \centering
  \begin{tabular}{lll}
    \toprule
    View & Components & Dim. \\
    \midrule
    Actor (execution)  & full self (25) + full teammate (25)                & \\
                       & + 2 $\times$ coarse opponent (rel.\ pos.\ + rank, 4) & \\
                       & + team indicator (1)                                & \textbf{59} \\
    \midrule
    Critic (training)  & full state of all 4 karts ($4\times25$) + race meta (2) & \textbf{102} \\
    \bottomrule
  \end{tabular}
\end{table}

\subsection{Action space}

The policy outputs one of $21$ discrete controller actions combining steering
(five stick positions), acceleration/brake, drift, trick, and \emph{item use}.
A key environment-verification finding (Section~\ref{sec:verify}) is that the
item is fired by the \textbf{L} button on the GameCube controller; an earlier
mapping used the \texttt{X} button (look-behind), under which karts pick up items
but can never use them---fatal for a study about item timing.

\subsection{Rank-based team reward}

Let $N$ be the live field size and $p_i$ agent $i$'s race position
($1=$ first). Define the normalized rank value
\begin{equation}
  v_i \;=\; \frac{N - p_i}{N - 1} \;\in\; [0,1], \qquad v_i = 1 \text{ at first place.}
\end{equation}
The individual reward is the potential-based shaping of $v_i$, and the team
reward shares the team's mean rank value among both members
($\mathcal{T}(i)$ is agent $i$'s team):
\begin{align}
  r^{\text{ind}}_{i,t}  &= \gamma\,v_{i,t+1} - v_{i,t}, \\
  r^{\text{team}}_{i,t} &= \gamma\,\bar v_{\mathcal{T}(i),t+1} - \bar v_{\mathcal{T}(i),t},
  \qquad \bar v_{\mathcal{T}} = \tfrac{1}{|\mathcal{T}|}\textstyle\sum_{j\in\mathcal{T}} v_j .
\end{align}
A sparse \textsc{FinishReward} pays a bonus scaled by finishing rank, a
\textsc{LapReward} marks lap completion, and light \textsc{Speed} and
\textsc{MiniTurbo} terms aid early learnability. The total reward is a weighted
sum; the individual/team weights directly instantiate the proposal's
``balance individual rank with team-level success.'' Because $r^{\text{ind}}$
and $r^{\text{team}}$ are potential-based, this dense shaping does not change the
optimal policy~\citep{shaping}.

\subsection{MAPPO with a centralized critic}

A single shared actor $\pi_\theta(a\mid o)$ conditions on the $59$-d restricted
observation $o$; a centralized critic $V_\phi(s)$ conditions on the $102$-d
global state $s$. Advantages are computed by GAE~\citep{gae} using the
centralized value, and the actor is updated with the clipped PPO
objective~\citep{ppo}:
\begin{equation}
  \mathcal{L}^{\text{actor}}(\theta) =
  \mathbb{E}\!\left[\min\!\big(\rho_t \hat A_t,\;
  \mathrm{clip}(\rho_t, 1-\epsilon, 1+\epsilon)\hat A_t\big)\right]
  + \beta\, \mathcal{H}[\pi_\theta], \qquad
  \rho_t = \frac{\pi_\theta(a_t\mid o_t)}{\pi_{\theta_{\text{old}}}(a_t\mid o_t)},
\end{equation}
with the critic trained by mean-squared error to the GAE returns. Centralization
is enforced architecturally: the critic only accepts the $102$-d state, so it
cannot silently degrade to a decentralized value function.

\subsection{Self-play}

Each episode, the learner team $\{0,1\}$ is driven by the live shared policy and
the opponent team $\{2,3\}$ by a \emph{frozen} snapshot sampled from a bounded
pool; only the learner team's transitions are used for the MAPPO update. The
current policy is periodically pushed into the pool. The win rate of the live
policy against sampled snapshots is logged as the proposal's training-health
signal: a rising win-rate-vs-older-snapshots indicates genuine improvement rather
than cycling.

\subsection{Kart symmetry and training curriculum}
\label{sec:controls}

\paragraph{Identical karts (RQ1 control).} For RQ1 to attribute behavioral
differences to \emph{rank} rather than to vehicle statistics, all four karts must
be identical. Mario Kart Wii assigns small per-character stat bonuses, so even
same-weight characters differ slightly; we therefore fix every agent to the
\emph{same} character and vehicle (Funky Kong on the Flame Runner bike). The
in-game character select is performed by a scripted menu macro, which we found
does not reliably place all four player cursors on the same cell---reading the
selected characters back from game memory showed only one of four matching the
target. We therefore initialize every episode by loading a verified save state in
which all four karts are confirmed identical (Funky Kong, character id $22$),
guaranteeing symmetry independently of menu navigation.

\paragraph{Race-first curriculum.} Learning to race, to use items, and to
cooperate simultaneously is difficult, and self-play under the full task is
observed to plateau at partial track completion. We therefore split training into
two phases that share an action space (and hence transfer by warm-starting). In
the \emph{racing} phase the item button is disabled for the learning agents, so
the policy concentrates on the racing line and vehicle control without
item-management noise; the learned weights then warm-start the \emph{full} phase,
in which items are re-enabled for the rank-conditioned study. We additionally vary
the episode start state across points further along the track (a diverse-start
curriculum), giving the agent concentrated practice on later sections it would
otherwise rarely reach. Both curricula are implemented as configuration switches
over the same network, so a policy trained under one transfers directly to the
other.


\section{Implementation and Verification}
\label{sec:verify}

The contributions are implemented as a self-contained package that reuses the
emulator/environment as a library. Each component is covered by an offline unit
test (no emulator required) and the full stack has been run end-to-end against
the live environment.

\paragraph{Environment verification (gating findings).} Before training we
verified the environment with a scripted probe.
(1)~\textbf{Items are enabled and are awarded}: karts pick up real item slots
while racing. (2)~\textbf{Item firing requires the \texttt{L} button}: with
\texttt{L}, held items discharge (confirmed on multiple karts); with the
previously used \texttt{X} button, zero items ever fired. We corrected the action
mapping accordingly. (3)~The default race is a \textbf{12-kart field} (four
agents plus eight CPUs); Mario Kart Wii's VS mode always spawns twelve karts, so
a literal ``four-kart'' race is not attainable. Rather than removing opponents, we
remove the \emph{confound} they would introduce into the rank analysis by giving
all four agents an identical kart (Section~\ref{sec:controls}), so any
rank-conditioned behavioral difference is attributable to rank rather than to
vehicle statistics. (4)~The memory reader is occasionally not ready immediately after a
race loads; our rollout wraps environment resets in a retry-with-warmup loop.

\paragraph{Unit and integration tests.} Table~\ref{tab:status} lists the verified
components. The asymmetry test confirms that an opponent's item is absent from
the actor observation but present in the critic state. The MAPPO test confirms
that a single update optimizes both networks and that the critic rejects a
restricted observation (i.e.\ it is genuinely centralized). The reward test
confirms overtakes are rewarded, a teammate's gain is shared identically by both
team members, and the finish bonus scales with rank and fires once. The
self-play test confirms snapshots are frozen and independent of the live policy.

\begin{table}[t]
  \caption{Implementation status and verification. All offline tests pass; the
  MAPPO stack additionally ran end-to-end against the live emulator.}
  \label{tab:status}
  \centering
  \begin{tabular}{lll}
    \toprule
    Component & Realization & Verification \\
    \midrule
    Asymmetric obs (CTDE)   & actor 59-d / critic 102-d            & opponent item hidden from actor \\
    Action space            & 21 discrete, item on \texttt{L}      & items fire end-to-end \\
    Rank--team reward       & potential-based, individual+team     & overtake / shared-team / finish \\
    MAPPO learner           & shared actor + centralized critic    & both nets update; critic centralized \\
    Self-play pool          & bounded frozen-snapshot pool         & snapshots frozen \& independent \\
    Training loop           & rollout + update + checkpoints       & ran end-to-end; clean shutdown \\
    \bottomrule
  \end{tabular}
\end{table}

\paragraph{End-to-end run.} The integrated trainer was executed against the live
four-player environment: observations and global state are built from real game
memory, the shared actor acts, the centralized critic and actor are updated by
MAPPO, and checkpoints are written. A graceful-shutdown handler saves a final
checkpoint and closes the experiment cleanly on \texttt{SIGTERM}, so interrupted
runs are recorded as finished rather than crashed.

\paragraph{Preliminary training observations.} End-to-end self-play produces
visibly competent driving and a win rate against frozen snapshots in the
neighborhood of $70\%$, evidence of genuine policy improvement rather than
strategy cycling. Track completion, however, plateaus short of finishing the full
race under the unrestricted task; this is what motivates the race-first and
diverse-start curricula of Section~\ref{sec:controls}. Quantitative learning
curves and the RQ1/RQ2 analyses are in progress at the time of writing and are
logged to a tracking dashboard for inclusion in the final version.


\section{Planned Experiments and Analysis}
\label{sec:plan}

With the verified system, the remaining work is the empirical study, on a fixed
track in the 2v2 setting with four identical karts (Section~\ref{sec:controls})
and items enabled.

\paragraph{RQ1: rank-differentiated behavior.} We will measure the divergence of
rank-conditioned action distributions, principally the mean number of frames an
item is \emph{held} before use, conditioned on the holder's rank. Strong
separation between rank~1 (long holding / defensive) and rank~4 (immediate use /
aggressive) indicates differentiation; near-identical distributions indicate mode
collapse. We will report item-holding-time and item-use-rate as a function of
rank.

\paragraph{RQ2: team cooperation.} We will count teammate-protective actions and
opponent-targeted attacks, summarized as the ratio of opponent hits to
teammate hits (``teamkills''). Detecting hit events requires a spin-out/collision
signal not currently read from memory; we will either add the corresponding
address or use a proxy (an opponent's abrupt speed drop coincident with our item
use).

\paragraph{Training health.} Throughout training we will track the win rate of
the current policy against frozen past snapshots; a rising curve evidences genuine
improvement. We will additionally report whether domain behaviors (mini-turbo
usage, item holding time, overtaking counts) align with documented human
strategies.


\section{Discussion and Limitations}

Two practical constraints shape the experimental plan. First, the environment
runs a single emulator instance and is stepped frame-by-frame, so wall-clock
throughput is modest; parallel instances would multiply it roughly linearly.
Second, because VS mode always fields twelve karts, the clean ``rank~1 vs.\
rank~4'' contrast is obtained not by removing opponents but by making the four
agents' karts identical (Section~\ref{sec:controls}), isolating rank from
vehicle-stat confounds. The hit/teamkill metric for RQ2 is the most uncertain
instrumentation and may rely on the speed-drop proxy. None of these
affect the implemented learning system, which is complete and verified.


\section{Conclusion}

We have specified and implemented a complete 2v2 MAPPO self-play system for Mario
Kart Wii: information-asymmetric CTDE observations, a potential-based rank reward
balancing individual and team success, a centralized-critic MAPPO learner, and a
frozen-snapshot self-play loop, together with a corrected item-control mapping
discovered during environment verification. Each component is unit-tested and the
full stack runs end-to-end. The remaining work is the empirical study of
rank-differentiated item strategy (RQ1) and team cooperation (RQ2) under the
protocol above.


\begin{thebibliography}{9}
\small

\bibitem{btr}
Clark, et al. (2025) Beyond the Rainbow: High-Performance Deep Reinforcement
Learning on a Desktop PC. \textit{International Conference on Machine Learning
(ICML)}.

\bibitem{mappo}
Yu, C., Velu, A., Vinitsky, E., Gao, J., Wang, Y., Bayen, A., \& Wu, Y. (2022)
The Surprising Effectiveness of PPO in Cooperative Multi-Agent Games.
\textit{Advances in Neural Information Processing Systems (NeurIPS)}.

\bibitem{ppo}
Schulman, J., Wolski, F., Dhariwal, P., Radford, A., \& Klimov, O. (2017)
Proximal Policy Optimization Algorithms. \textit{arXiv:1707.06347}.

\bibitem{maddpg}
Lowe, R., Wu, Y., Tamar, A., Harb, J., Abbeel, P., \& Mordatch, I. (2017)
Multi-Agent Actor-Critic for Mixed Cooperative-Competitive Environments.
\textit{Advances in Neural Information Processing Systems (NeurIPS)}.

\bibitem{shaping}
Ng, A.\,Y., Harada, D., \& Russell, S. (1999) Policy Invariance Under Reward
Transformations: Theory and Application to Reward Shaping.
\textit{International Conference on Machine Learning (ICML)}.

\bibitem{gae}
Schulman, J., Moritz, P., Levine, S., Jordan, M., \& Abbeel, P. (2016)
High-Dimensional Continuous Control Using Generalized Advantage Estimation.
\textit{International Conference on Learning Representations (ICLR)}.

\bibitem{selfplay}
Bansal, T., Pachocki, J., Sidor, S., Sutskever, I., \& Mordatch, I. (2018)
Emergent Complexity via Multi-Agent Competition.
\textit{International Conference on Learning Representations (ICLR)}.

\end{thebibliography}


\end{document}
